% GENERATED FILE. Edit canonical lessons, then run npm run generate:books.
% canonical-source-hash: fnv1a64:d6c182905f306d2b
% canonical-lessons: PT-C451-sapateiro, PT-C451-costureira, PT-C451-bainha, PT-C451-serralheiro, PT-C451-remendar

\chapter{Arranjar e remendar}
\label{ch:pt-arranjar-e-remendar}
\hlchaptermodality{451}

\begin{chapteropening}
\textbf{By the end of this chapter:} \emph{I can get things mended: the cobbler, the seamstress, the hem, the locksmith and patching.}

Complete the last lesson of chapter 451: I can get things mended: the cobbler, the seamstress, the hem, the locksmith and patching.
\end{chapteropening}

\section[o sapateiro]{o sapateiro — the cobbler, the shoemaker}

\label{lesson:PT-C451-sapateiro}



\begin{quote}
Before the new one: say the Portuguese for the wine shop, then the Portuguese for to slice.
\end{quote}

\subsection*{You'll want to know: o sapateiro}
\textbf{o sapateiro} — \textquotedblleft{}the cobbler, the shoemaker\textquotedblright{}.

\begin{quote}
O sapateiro mudou as solas das minhas botas.
\end{quote}

\begin{quote}
\emph{The cobbler replaced the soles of my boots.}
\end{quote}

Built on \emph{o sapato}.

\begin{cousinweb}
\emph{sapateiro} is built on \emph{sapato} 'shoe' with \emph{-eiro}. \emph{Sapato} comes from Old Portuguese \emph{çapato} (EtymDB); its earlier origin is not settled.
\end{cousinweb}

\begin{grammarlens}[title={What you've built}]
One more word for telling, explaining and describing at B1.
\end{grammarlens}

\subsection*{Your turn}
\begin{itemize}
\raggedright
  \item \emph{Say it:} \emph{o sapateiro}
  \item \emph{Say it:} \emph{o sapateiro}, once more
  \item \emph{Say it:} a sentence of your own with \emph{o sapateiro}
  \item \emph{Recall:} say the Portuguese for the wine shop, then the Portuguese for to slice, then say \emph{o sapateiro} again
\end{itemize}

\begin{tcolorbox}[breakable,colback=teal!4,colframe=teal!35!black,title={Before you move on}]
What does o sapateiro mean? (\textquotedblleft{}The cobbler, the shoemaker\textquotedblright{}.) Say it once more.
\end{tcolorbox}
\section[a costureira]{a costureira — the seamstress, the dressmaker}

\label{lesson:PT-C451-costureira}



\begin{quote}
Before the new one: say the Portuguese for to slice, then the Portuguese for the cobbler.
\end{quote}

\subsection*{You'll want to know: a costureira}
\textbf{a costureira} — \textquotedblleft{}the seamstress, the dressmaker\textquotedblright{}.

\begin{quote}
A costureira arranjou o meu vestido.
\end{quote}

\begin{quote}
\emph{The seamstress altered my dress.}
\end{quote}

Masculine: \emph{o costureiro}, often a fashion designer.

\begin{cousinweb}
\emph{costureira} is built on \emph{costura} 'sewing' with \emph{-eira}, the person who does it; \emph{costura} goes back to Vulgar Latin \emph{consūtūra}, from Latin \emph{consuere} 'to sew together', on \emph{suere} 'to sew' (Lewis \& Short).
\end{cousinweb}

\begin{grammarlens}[title={What you've built}]
One more word for telling, explaining and describing at B1.
\end{grammarlens}

\subsection*{Your turn}
\begin{itemize}
\raggedright
  \item \emph{Say it:} \emph{a costureira}
  \item \emph{Say it:} \emph{a costureira}, once more
  \item \emph{Say it:} a sentence of your own with \emph{a costureira}
  \item \emph{Recall:} say the Portuguese for to slice, then the Portuguese for the cobbler, then say \emph{a costureira} again
\end{itemize}

\begin{tcolorbox}[breakable,colback=teal!4,colframe=teal!35!black,title={Before you move on}]
What does a costureira mean? (\textquotedblleft{}The seamstress, the dressmaker\textquotedblright{}.) Say it once more.
\end{tcolorbox}
\section[a bainha]{a bainha — the hem; the sheath}

\label{lesson:PT-C451-bainha}



\begin{quote}
Before the new one: say the Portuguese for the cobbler, then the Portuguese for the seamstress.
\end{quote}

\subsection*{You'll want to know: a bainha}
\textbf{a bainha} — \textquotedblleft{}the hem; the sheath\textquotedblright{}.

\begin{quote}
A costureira subiu a bainha das calças.
\end{quote}

\begin{quote}
\emph{The seamstress took up the hem of the trousers.}
\end{quote}

\emph{Subir a bainha} — to take up a hem.

\begin{cousinweb}
\emph{bainha} is inherited from Latin \emph{vāgīna} 'sheath, scabbard' (EtymDB): the folded edge holds the cloth as a sheath holds a sword.
\end{cousinweb}

\begin{grammarlens}[title={What you've built}]
One more word for telling, explaining and describing at B1.
\end{grammarlens}

\subsection*{Your turn}
\begin{itemize}
\raggedright
  \item \emph{Say it:} \emph{a bainha}
  \item \emph{Say it:} \emph{a bainha}, once more
  \item \emph{Say it:} a sentence of your own with \emph{a bainha}
  \item \emph{Recall:} say the Portuguese for the cobbler, then the Portuguese for the seamstress, then say \emph{a bainha} again
\end{itemize}

\begin{tcolorbox}[breakable,colback=teal!4,colframe=teal!35!black,title={Before you move on}]
What does a bainha mean? (\textquotedblleft{}The hem; the sheath\textquotedblright{}.) Say it once more.
\end{tcolorbox}
\section[o serralheiro]{o serralheiro — the locksmith; the metalworker}

\label{lesson:PT-C451-serralheiro}



\begin{quote}
Before the new one: say the Portuguese for the seamstress, then the Portuguese for the hem.
\end{quote}

\subsection*{You'll want to know: o serralheiro}
\textbf{o serralheiro} — \textquotedblleft{}the locksmith; the metalworker\textquotedblright{}.

\begin{quote}
Perdi a chave e chamei o serralheiro.
\end{quote}

\begin{quote}
\emph{I lost my key and called the locksmith.}
\end{quote}

He makes gates and railings too.

\begin{cousinweb}
\emph{serralheiro} is usually traced through Old Portuguese \emph{serralha} 'lock' to Latin \emph{sera} 'bar, bolt' (Lewis \& Short).
\end{cousinweb}

\begin{grammarlens}[title={What you've built}]
One more word for telling, explaining and describing at B1.
\end{grammarlens}

\subsection*{Your turn}
\begin{itemize}
\raggedright
  \item \emph{Say it:} \emph{o serralheiro}
  \item \emph{Say it:} \emph{o serralheiro}, once more
  \item \emph{Say it:} a sentence of your own with \emph{o serralheiro}
  \item \emph{Recall:} say the Portuguese for the seamstress, then the Portuguese for the hem, then say \emph{o serralheiro} again
\end{itemize}

\begin{tcolorbox}[breakable,colback=teal!4,colframe=teal!35!black,title={Before you move on}]
What does o serralheiro mean? (\textquotedblleft{}The locksmith; the metalworker\textquotedblright{}.) Say it once more.
\end{tcolorbox}
\section[remendar]{remendar — to patch, to mend}

\label{lesson:PT-C451-remendar}



\begin{quote}
Before the new one: say the Portuguese for the hem, then the Portuguese for the locksmith.
\end{quote}

\subsection*{You'll want to know: remendar}
\textbf{remendar} — \textquotedblleft{}to patch, to mend\textquotedblright{}.

\begin{quote}
A minha avó remenda as meias à noite.
\end{quote}

\begin{quote}
\emph{My grandmother patches socks in the evening.}
\end{quote}

\emph{Um remendo} — a patch.

\begin{cousinweb}
\emph{remendar} is built from \emph{re-} and \emph{emendar} 'to correct' (EtymDB), which is inherited from Latin \emph{ēmendāre} 'to free from faults'.
\end{cousinweb}

\begin{grammarlens}[title={What you've built}]
One more word for telling, explaining and describing at B1.
\end{grammarlens}

\subsection*{Your turn}
\begin{itemize}
\raggedright
  \item \emph{Say it:} \emph{remendar}
  \item \emph{Say it:} \emph{remendar}, once more
  \item \emph{Say it:} a sentence of your own with \emph{remendar}
  \item \emph{Recall:} say the Portuguese for the hem, then the Portuguese for the locksmith, then say \emph{remendar} again
\end{itemize}

\begin{tcolorbox}[breakable,colback=teal!4,colframe=teal!35!black,title={Before you move on}]
What does remendar mean? (\textquotedblleft{}To patch, to mend\textquotedblright{}.) Say it once more.
\end{tcolorbox}
